\documentclass[11pt]{article}

\usepackage[utf8]{inputenc}
\usepackage[T1]{fontenc}
\usepackage{lmodern}
\usepackage{amsmath,amssymb}
\usepackage{graphicx}
\usepackage{booktabs}
\usepackage{microtype}
\usepackage{geometry}
\usepackage{caption}
\usepackage{natbib}
\usepackage{hyperref}

\geometry{margin=1in}
\setcitestyle{numbers,square}

\title{Antipsychotic drugs selectively decorrelate long-range cortical interactions mediated by layer 5 intratelencephalic neurons}
\author{}
\date{}

\begin{document}
\maketitle

\begin{abstract}
Antipsychotic drugs are effective against positive symptoms of psychotic disorders, yet how their action reshapes distributed cortical circuits has remained obscure. Existing pharmacological accounts emphasise dopamine D2 and 5-HT2A antagonism, but do not specify where in the cortex---or in which cell classes---clinical efficacy is reflected in coordinated activity. To address this gap, we combined widefield calcium imaging across twelve dorsal cortical regions with cell-type-specific expression of GCaMP6 in head-fixed mice navigating a virtual corridor. Pan-neuronal expression was achieved by retro-orbital AAV-PHP.eB in C57BL/6 mice; layer 5 intratelencephalic (L5 IT) neurons and layer 2/3 neurons were targeted using Tlx3-Cre $\times$ Ai148 and Cux2-CreERT2 $\times$ Ai148 crosses, respectively. Imaging was performed at 100 Hz through the cleared skull before and after single intraperitoneal injections of clozapine, aripiprazole, haloperidol, amphetamine, or saline. Pearson correlations between pairs of regions were split at $0.9\times$ the bregma--lambda distance (${\sim}3.8$~mm) into short- and long-range pairs. Antipsychotic treatment produced a significantly stronger reduction of both short- and long-range correlations in L5 IT neurons than in pan-neuronal recordings (C57BL/6 vs Tlx3-Cre $\times$ Ai148: short-range $p<0.05$; long-range $p<10^{-5}$, rank-sum test), and the reduction was significantly weaker in L2/3 Cux2 neurons than in L5 IT Tlx3 neurons (short-range $p<0.005$; long-range $p<10^{-8}$). The psychostimulant amphetamine did not reproduce the antipsychotic-induced decorrelation. These results identify L5 IT pyramidal neurons as a cell-type-specific cortical substrate of antipsychotic action and provide a circuit-level framework linking receptor pharmacology to distributed cortical coupling.
\end{abstract}

\section{Introduction}

The cerebral cortex supports perception, action, and internal-model inference by coordinating activity across many anatomically distributed areas. This long-range coordination is thought to depend on specific classes of pyramidal neurons whose axonal projections span the cortical mantle, and its disruption has repeatedly been invoked as a substrate of psychotic disorders~\citep{uhlhaas2015oscillations, uhlhaas2013highfrequency, bahner2017hippocampal}. Hierarchical, anatomically graded connectivity between cortical areas provides the scaffold on which this coordination is built~\citep{harris2019hierarchical, siegle2021survey}, and different classes of glutamatergic neurons contribute differently to inter-areal and cortico-cortical signalling~\citep{matho2021genetic, biccn2021multimodal, liu2024celltype}. Despite decades of progress on the receptor pharmacology of antipsychotic drugs, how their therapeutic action reshapes this distributed cortical circuit remains poorly understood.

The dominant pharmacological account of antipsychotic efficacy emphasises dopamine D2 receptor antagonism, with complementary actions at 5-HT2A and other monoaminergic targets~\citep{howes2009dopamine, kambeitz2014alterations}. While this framework has guided drug development for half a century, it does not specify where in the cortex antipsychotics change neural activity, nor whether their effect is uniform across cortical layers and cell classes. Because positive symptoms of psychosis are associated with aberrant salience and disordered long-range synchrony~\citep{uhlhaas2015oscillations, bahner2017hippocampal}, a cortical circuit account of antipsychotic action must go beyond receptor occupancy and ask how receptor binding translates into altered inter-areal coupling. The key gap we address here is a cell-type-resolved, within-mouse measurement of inter-areal correlation before and after antipsychotic administration.

Layer 5 intratelencephalic (L5 IT) pyramidal neurons are attractive candidate substrates for such coupling. They project densely within cortex and between hemispheres, and recent work using genetic labelling has delineated their anatomical reach and functional contributions to cortico-cortical signalling~\citep{liu2024celltype, shinotsuka2023layer5, atudorei2024bilateral, matho2021genetic}. The Tlx3-Cre driver line targets L5 IT neurons selectively; combined with Cre-dependent reporter lines such as Ai148 that express GCaMP6 under Cre control~\citep{daigle2018suite, chen2013ultrasensitive}, this driver enables cell-type- and layer-specific readouts of distributed cortical activity. The Cux2-CreERT2 line, combined with tamoxifen-induced recombination, provides comparable access to layer 2/3 pyramidal neurons and thus a crucial layer-control comparison~\citep{daigle2018suite}. Pan-neuronal expression, in turn, can be achieved by the retro-orbitally administered AAV-PHP.eB capsid~\citep{chan2017engineered}, yielding a third, genetically unselective population against which the cell-type-specific measurements can be benchmarked.

Widefield calcium imaging across the dorsal cortex of head-fixed mice has emerged as a powerful tool for capturing distributed activity at a scale commensurate with inter-areal coupling~\citep{musall2019single, leinweber2017sensorimotor}. Combined with a head-fixed virtual reality setup, it permits the alignment of inter-areal activity to behaviourally well-defined events such as closed- and open-loop locomotion onsets and visual-flow onsets~\citep{leinweber2017sensorimotor, keller2012sensorimotor, attinger2017visuomotor, keller2018predictive}. Under a predictive-processing framing, these events separate bottom-up and top-down contributions to cortical activity and thus provide a behaviourally anchored readout through which drug-induced changes in coupling can be interpreted~\citep{keller2018predictive, zmarz2016mismatch}. Methodological refinements---including hierarchical bootstrap statistics for nested neuronal data~\citep{saravanan2020hierarchical} and percentile-based baseline correction adapted from two-photon imaging~\citep{dombeck2007imaging}---allow this approach to be combined with pharmacological perturbations within individual mice, so that the same animal serves as its own control.

We combine these tools to ask whether antipsychotic drugs selectively alter long-range cortical coupling, and whether this alteration is specific to a molecularly defined subclass of pyramidal neurons. We image twelve regions across dorsal cortex before and after single intraperitoneal injections of clozapine, aripiprazole, haloperidol, amphetamine, or saline in each of the three populations above, and quantify Pearson correlations between pairs of regions as a function of inter-region distance. Our results make three complementary contributions. First, antipsychotics reduce short- and long-range correlations in L5 IT neurons substantially more than they do in pan-neuronal recordings. Second, the effect is substantially weaker in L2/3 Cux2 neurons than in L5 IT Tlx3 neurons, indicating layer specificity. Third, the psychostimulant amphetamine, which shares dopaminergic but not antipsychotic properties, does not reproduce the decorrelation pattern. Together, these observations identify L5 IT neurons as a cell-type-specific cortical substrate of antipsychotic action and support a circuit-level account of therapeutic efficacy that complements receptor-centric models.

\section{Background}

\paragraph{Cortical interactions and psychosis.}
The idea that psychosis involves altered long-range cortical coordination is supported by both electrophysiological and hemodynamic evidence. Abnormal beta- and gamma-band oscillations, and reduced large-scale synchrony, have been repeatedly reported in schizophrenia~\citep{uhlhaas2015oscillations, uhlhaas2013highfrequency}. Hippocampal--prefrontal coupling during cognitive tasks has been proposed as a translational phenotype whose alterations track illness and genetic risk~\citep{bahner2017hippocampal}. These observations motivate measurements of inter-areal correlation in animal models of pharmacological intervention and justify treating decorrelation, rather than mean firing rate, as a primary outcome.

\paragraph{Dopaminergic and receptor-level accounts.}
The dopamine hypothesis of schizophrenia, and the aberrant-salience framework that grew out of it, anchor antipsychotic action in D2 receptor antagonism and associated dopamine dysfunction~\citep{howes2009dopamine, kambeitz2014alterations}. Clozapine, aripiprazole, and haloperidol differ substantially in their affinities for D2, 5-HT2A, and other targets, yet are all clinically effective against positive symptoms. This clinical convergence despite pharmacological divergence suggests a shared downstream effect at the circuit level that receptor-centric models do not fully specify, and which our experiments are designed to expose.

\paragraph{Genetic and anatomical dissection of cortical cell types.}
The past two decades have produced Cre driver lines and anatomical atlases that enable cell-type-specific access to cortical pyramidal neuron subclasses. L5 IT neurons, labelled by Tlx3-Cre, form dense cortico-cortical and cortico-striatal projections~\citep{liu2024celltype, matho2021genetic, biccn2021multimodal, shinotsuka2023layer5, atudorei2024bilateral}, and their activity contributes to coordinated behaviours that require distributed cortical signals. Layer 2/3 neurons in sensory cortex, by contrast, participate in predominantly local computations and are targeted by the Cux2-CreERT2 line~\citep{daigle2018suite}. Hierarchically organised cortico-cortical and cortico-thalamic connectivity provides the structural backbone for the interactions measured here~\citep{harris2019hierarchical, siegle2021survey}.

\paragraph{Widefield imaging and cell-type tools.}
Widefield imaging of genetically encoded calcium indicators such as GCaMP6~\citep{chen2013ultrasensitive} captures distributed cortical activity at a spatial scale that is well matched to the footprint of inter-areal coupling~\citep{musall2019single, leinweber2017sensorimotor}. Cre-dependent transgenic reporter lines including Ai148~\citep{daigle2018suite, bounds2023alloptical} make cell-type-specific widefield recordings practical, and retro-orbital delivery of AAV-PHP.eB enables pan-neuronal or DIO-gated expression without intracranial injection~\citep{chan2017engineered}. Together, these tools allow the same imaging pipeline to be applied across genetically defined populations, providing the controlled cell-type comparison on which the present study rests.

\paragraph{Predictive processing and sensorimotor integration.}
A complementary line of work has framed cortical activity during locomotion and visual flow in terms of predictive processing~\citep{keller2012sensorimotor, keller2018predictive, zmarz2016mismatch, attinger2017visuomotor, leinweber2017sensorimotor}. Closed-loop and open-loop locomotion onsets, as well as mismatches between predicted and actual visual flow, reveal separable bottom-up and top-down contributions to cortical activity. Top-down projections from frontal and cingulate areas shape activity in primary visual cortex and carry signals consistent with predictions of visual flow based on motor output~\citep{leinweber2017sensorimotor}. These behaviourally defined events serve as internal benchmarks against which we compare naive and drug-treated activity.

\section{Methods}

\subsection{Animals and genetic strategies}
All experiments used adult mice of either sex. Nine mouse lines were employed: C57BL/6, Emx1-Cre, Scnn1a-Cre, Tlx3-Cre~$\times$~Ai148, Ntsr1-Cre, PV-Cre, VIP-Cre, Sst-Cre, and Cux2-CreERT2~$\times$~Ai148. To induce Cre activity in the Cux2-CreERT2 line, mice were fed tamoxifen-containing food (Envigo, 400~mg tamoxifen per kg of food) as their sole food source for at least one week prior to imaging. Experimental mouse counts per genotype (Figure~\ref{fig:mouse_counts}) were 6 C57BL/6, 4 Emx1-Cre, 7 Scnn1a-Cre, 14 Tlx3-Cre, 3 Ntsr1-Cre, 2 PV-Cre, 6 VIP-Cre, 6 Sst-Cre, and 3 Cux2-CreERT2 mice. Cell population, indicator, and delivery strategy for each cohort are summarised in Table~\ref{tab:cohorts}.

\subsection{Surgery and viral delivery}
Mice were anaesthetised using a mixture of fentanyl (0.05~mg/kg; Actavis), midazolam (5.0~mg/kg; Dormicum, Roche), and medetomidine (0.5~mg/kg; Domitor, Orion). Following cranial window implantation, anaesthesia was antagonised by an intraperitoneal injection of flumazenil (0.5~mg/kg; Anexate, Roche) and atipamezole (2.5~mg/kg; Antisedan, Orion Pharma). Peri- and post-operative analgesia used buprenorphine (0.1~mg/kg; Reckitt Benckiser) and Metacam (5~mg/kg; Boehringer Ingelheim). In a subset of mice (6 C57BL/6, 4 Emx1-Cre, and 4 Sst-Cre), an AAV vector based on the PHP.eB capsid~\citep{chan2017engineered} was injected retro-orbitally (6~$\mu$L per eye at $\geq 10^{13}$~GC/ml) to drive GCaMP expression under either EF1$\alpha$ or hSyn1 promoters pan-neuronally, or in a cell-type-specific manner using a double-floxed inverted open reading frame (DIO) construct. Imaging commenced at the earliest 1 week after head-bar implantation or 3 weeks after retro-orbital AAV injection.

\subsection{Pharmacological agents}
Clozapine (Novartis), aripiprazole (Otsuka Pharmaceutical), haloperidol (Janssen), amphetamine (H\"anseler), or saline (0.9\% NaCl in H$_2$O) was administered intraperitoneally at doses of 0.2--10~$\mu$g/g, 0.2~$\mu$g/g, 0.1~$\mu$g/g, 4~$\mu$g/g body weight, or undiluted, respectively. Each mouse received a single injection, and within-session data were aggregated across pre-injection (``naive'') and post-injection epochs (see Table~\ref{tab:drugs}).

\subsection{Virtual reality setup}
All experiments used the virtual reality setup previously described by~\citet{leinweber2017sensorimotor}, in which head-fixed mice ran on a spherical treadmill to navigate a virtual corridor. Closed-loop locomotion onsets coupled the treadmill motion to visual flow; open-loop locomotion onsets and isolated visual-flow onsets allowed dissociation of motor and visual contributions. Well-isolated onsets were selected with a 30~cm/s speed threshold, and we excluded all onsets with locomotion or visual flow during the 3~s window preceding the onset. The baseline subtraction window for unpredictable stimuli (mismatches and grating onsets) was $-200$~ms to 0~ms before stimulus onset; for closed- and open-loop locomotion onsets and open-loop visual-flow onsets it was $-2900$~ms to $-2700$~ms for GCaMP imaging or $-900$~ms to $-700$~ms for eGFP imaging, chosen to accommodate calcium indicator offset dynamics and potential anticipatory activity.

\subsection{Widefield macroscope}
Widefield imaging was performed on a custom-built macroscope consisting of two commercially available objectives mounted face-to-face (Nikon 85~mm/$f$1.8 sample side; Nikon 50~mm/$f$1.4 sensor side). Fluorescence was excited by a 470~nm LED (Thorlabs) driven at a 90\% duty cycle, passed through an SP490 excitation filter (Thorlabs), reflected off an LP490 dichroic mirror (Thorlabs), and collected through a 525/50 emission filter onto a PCO edge 4.2 sCMOS camera. Apertures on the objectives were generally kept maximally open, and the LED driver current was adjusted to keep fluorescence below 25\% of the sensor's maximum dynamic range. Raw images were acquired at 100~Hz (except for eGFP controls at 50~Hz effective frame rate) at 16-bit depth. Imaging covered twelve regions of interest (ROIs) on the dorsal cortex. The experimental pipeline is summarised schematically in Figure~\ref{fig:schematic}.

\subsection{Two-photon imaging}
For complementary two-photon measurements, a tunable femtosecond laser (Coherent Chameleon) tuned to 930~nm was used with a 12~kHz resonant scanner (Cambridge Technology). Photomultiplier signals were amplified (DHPCA-100, Femto), digitised (NI5772, National Instruments) at 800~MHz, and band-pass filtered around 80~MHz using a digital Fourier-transform filter implemented on an FPGA (NI5772). Images were acquired at 750~$\times$~400~pixels and 60~Hz, and a piezo-electric actuator (P-726, Physik Instrumente) moved a Nikon 16$\times$ 0.8~NA objective through four planes in 15~$\mu$m steps, yielding an effective frame rate of 15~Hz over a 350~$\times$~350~$\mu$m field of view.

\subsection{Signal processing}
Off-line processing and analysis used custom-written MATLAB (2021a) scripts. For each ROI, activity was computed as $\Delta F / F_0$, where $F_0$ is the median fluorescence of the recording (approximately 30{,}000 frames in a 5~min recording). $\Delta F / F_0$ was corrected for slow fluorescence drift caused by thermal brightening of the LED using 8th-percentile filtering with a 62.5~s moving window, following a procedure described previously for two-photon imaging by~\citet{dombeck2007imaging}. Sessions in which the average activity of all 12 ROIs remained above 30\% $\Delta F / F$ for more than 10~s were excluded; this criterion removed 10 of 3801 five-minute recording sessions (0.26\%; Figure~\ref{fig:exclusion}). Given that the signal attenuation half-depth in cortex is on the order of 100~$\mu$m, most widefield signals are expected to arise from the superficial ${\sim}100~\mu$m of tissue. Hemodynamic correction using 405~nm isosbestic illumination was attempted but could not recover the hemodynamic responses measured in eGFP controls; eGFP-expressing Tlx3-Cre mice (6 mice, injected retro-orbitally with AAV-PHP.eB-Ef1$\alpha$-DIO-eGFP) therefore provided the critical control for hemodynamic contamination.

\subsection{Correlation analysis}
For the correlation analyses, Pearson correlation coefficients were computed between all possible pairs of the 12 ROIs. Correlations were then plotted against inter-region distance and binned into a $40\times40$ grid, which was smoothed with a Gaussian filter to produce the heatmaps from which contour lines at 50\% of the maximum value were drawn. Pair distances were split into short- and long-range sets using a cutoff of $0.9\times$ the bregma--lambda distance (approximately 3.8~mm). Pair counts for each comparison (Figure~\ref{fig:pair_counts}) were: 204 short-range and 192 long-range pairs for pan-neuronal C57BL/6 mice, 182 and 148 for Tlx3-Cre~$\times$~Ai148 mice treated with clozapine, 114 and 84 for Tlx3-Cre~$\times$~Ai148 mice treated with haloperidol or aripiprazole, and 137 and 127 for Cux2-CreERT2~$\times$~Ai148 mice treated with clozapine.

\subsection{Statistics}
For widefield time-course data (Figures~1F--1L, 3B--3C, 3E--3F, 4, 8, and related supplements in the source dataset), we used a hierarchical bootstrap estimate of the mean and its 90\% confidence interval~\citep{saravanan2020hierarchical}, generated from 1000 bootstrap samples chosen hierarchically (first over mice, then over onsets) for each 10~ms time bin. Closed- versus open-loop comparisons used one-way ANOVA with genotype or mouse as factor, followed by a bootstrap-distribution-based estimate of $p$-values. Correlation comparisons (Figures~5--7 and related supplements) used rank-sum tests with family-wise error correction: the adjusted threshold was $\alpha_{\mathrm{adj}} = \alpha / m$ with $m = 9$ groups (genotypes), and significance indicators reflect $\alpha \in \{0.05, 0.01, 0.001\}$. Cross-genotype $p$-value bounds are summarised in Figure~\ref{fig:pvalues} and Table~\ref{tab:stats}. Mice in which the average locomotion-onset response was below 1\% $\Delta F / F$ in either the closed- or open-loop condition were excluded from the onset-similarity correlation analysis (3 of 54 mice), because the correlation analysis yields noise in the absence of a clear response.

\section{Results}

\subsection{Widefield imaging captures reproducible cortical activity during virtual navigation}

To measure inter-areal cortical coupling during naturalistic behaviour, we imaged widefield GCaMP fluorescence across twelve regions of dorsal cortex in head-fixed mice navigating a virtual corridor. Our custom macroscope delivered 470~nm LED illumination at a 90\% duty cycle and acquired 16-bit images at 100~Hz through a PCO edge 4.2 sCMOS camera (Figure~\ref{fig:schematic}). With these parameters, both closed-loop and open-loop locomotion onsets, and isolated visual-flow onsets in the open-loop condition restricted to periods when mice were not locomoting, evoked structured responses that averaged cleanly across trials. Mean responses were estimated with a hierarchical bootstrap (1000 samples over mice and then over onsets per 10~ms time bin) and reported with 90\% confidence intervals~\citep{saravanan2020hierarchical}.

To ensure data quality, we excluded sessions in which the average activity across all twelve ROIs remained above 30\% $\Delta F / F$ for more than 10~s. This criterion excluded 10 of 3801 five-minute recording sessions, or 0.26\% of the data (Figure~\ref{fig:exclusion}), indicating that the widefield recordings were generally well-behaved and that the exclusion criterion was conservative. Across genotypes, the similarity of closed- and open-loop locomotion onset responses was assessed by a one-way ANOVA with genotype as factor, followed by a bootstrap-based $p$-value estimate. Mice whose onset responses were below 1\% $\Delta F / F$ in either condition (3 of 54 mice) were excluded from this particular analysis, because the correlation between small, noisy responses yields uninformative estimates. Together, these checks confirmed that the widefield imaging pipeline provided a reproducible substrate for the cell-type and pharmacological comparisons that follow.

\subsection{Cell-type-specific GCaMP6 expression enables layer-resolved comparisons}

Having established a stable imaging pipeline, we next configured it for cell-type-specific readout. Three strategies delivered GCaMP6 to complementary populations. Pan-neuronal expression was achieved by retro-orbital injection of AAV-PHP.eB~\citep{chan2017engineered} into 6 C57BL/6 mice, 4 Emx1-Cre mice, and 4 Sst-Cre mice, driven by EF1$\alpha$ or hSyn1 promoters (or DIO constructs for cell-type-specific expression). L5 IT neurons were targeted via the Tlx3-Cre $\times$ Ai148 cross, which expresses GCaMP6 under Cre control in L5 IT pyramidal cells~\citep{daigle2018suite, liu2024celltype}. Layer 2/3 neurons were targeted via the Cux2-CreERT2 $\times$ Ai148 cross, with Cre activity induced by at least one week of tamoxifen-containing food (400~mg/kg). Mouse counts used for onset-response analyses across nine genotype groups are summarised in Figure~\ref{fig:mouse_counts} and in Table~\ref{tab:cohorts}: 6 C57BL/6, 4 Emx1-Cre, 7 Scnn1a-Cre, 14 Tlx3-Cre, 3 Ntsr1-Cre, 2 PV-Cre, 6 VIP-Cre, 6 Sst-Cre, and 3 Cux2-CreERT2 mice.

Because widefield fluorescence can in principle be contaminated by hemodynamic signals, we performed an eGFP control in 6 Tlx3-Cre mice injected retro-orbitally with AAV-PHP.eB-Ef1$\alpha$-DIO-eGFP. These mice were imaged at an effective frame rate of 50~Hz with otherwise identical acquisition parameters. We attempted to correct the GCaMP signal for hemodynamic response using a 405~nm isosbestic illumination, but could not fully recover the hemodynamic component measured in the eGFP control, consistent with the view that accurate correction requires calibration on eGFP measurements. The eGFP data therefore served as the principal control for hemodynamic contamination in the drug-effect analyses, rather than an online 405~nm correction. This control design ensured that any cell-type-specific decorrelation we subsequently observed in GCaMP recordings could not be trivially attributed to vascular rather than neural origin.

\subsection{Clozapine modestly reduces cortical correlations when GCaMP is expressed pan-neuronally}

With the imaging and cell-type strategies in place, we first asked how clozapine, the archetypal atypical antipsychotic, alters cortical correlations in the pan-neuronal condition. Four C57BL/6 mice expressing GCaMP6 via retro-orbital AAV-PHP.eB were imaged before and after a single intraperitoneal injection of clozapine, at doses spanning 0.2--10~$\mu$g/g body weight. For each state, Pearson correlations were computed across all pairs of the twelve ROIs and plotted against inter-region distance; the resulting heatmaps were binned on a $40\times40$ grid and smoothed with a Gaussian filter, with 50\% contour lines drawn on each panel and the pre-drug contour overlaid on the post-drug panel for comparison~\citep{dombeck2007imaging, saravanan2020hierarchical}.

At the pan-neuronal level, clozapine produced only a modest change in correlation structure. When pair-wise correlations were split into short- and long-range sets using the $0.9\times$ bregma--lambda cutoff ($\sim$3.8~mm), the pan-neuronal comparison drew on 204 short-range pairs and 192 long-range pairs (Figure~\ref{fig:pair_counts}; Table~\ref{tab:stats}). Statistical analysis of these pairs yielded significant changes in one range but not the other ($^{***}$: $p<0.001$; n.s.: not significant), indicating that, when GCaMP is distributed across all cortical neurons, the drug-induced decorrelation is partial and comparatively weak. We interpret this result as setting a conservative baseline for subsequent cell-type-specific comparisons: any stronger effect in a specific subclass must be resolved against this pan-neuronal reference.

\subsection{Antipsychotics strongly and selectively decorrelate L5 IT activity}

Having established a modest pan-neuronal baseline, we next targeted GCaMP6 expression to L5 IT neurons using the Tlx3-Cre $\times$ Ai148 cross. Five Tlx3-Cre $\times$ Ai148 mice were imaged before and after a single clozapine injection, using the same analysis pipeline. In contrast to the pan-neuronal result, clozapine produced significant reductions in both short-range (182 pairs) and long-range (148 pairs) correlations at the $^{***}$: $p<0.001$ level. Direct rank-sum comparison with the pan-neuronal baseline confirmed that the L5 IT reduction was significantly larger than the C57BL/6 reduction for both distance ranges: short-range $p<0.05$, long-range $p<10^{-5}$ (Figure~\ref{fig:pvalues}; Table~\ref{tab:stats}). Critically, the long-range effect size exceeded the short-range effect by several orders of magnitude in the bound on $p$, consistent with the idea that decorrelation primarily acts on distributed, inter-areal coupling rather than on local circuits.

To test whether the effect generalises beyond clozapine, we repeated the protocol with two further antipsychotics. Three Tlx3-Cre $\times$ Ai148 mice injected with aripiprazole (0.2~$\mu$g/g) showed a qualitatively similar decorrelation pattern in their heatmaps and contour lines. A separate group of three Tlx3-Cre $\times$ Ai148 mice treated with haloperidol (0.1~$\mu$g/g) yielded 114 short-range and 84 long-range pairs and produced a significant reduction in correlations (Figure~\ref{fig:pair_counts}; Table~\ref{tab:stats}), with short-range significance at $^{*}$: $p<0.05$ and long-range significance at $^{*}$: $p<0.05$ (n.s. where not significant) or $^{***}$: $p<0.001$ in the pooled analysis. Because clozapine, aripiprazole, and haloperidol differ substantially in their D2 and 5-HT2A profiles but converge on L5 IT decorrelation, the data suggest a shared downstream circuit effect that is not obviously predicted by receptor occupancy alone~\citep{howes2009dopamine, kambeitz2014alterations}.

\subsection{Decorrelation is weaker in L2/3 (Cux2) neurons than in L5 IT (Tlx3) neurons}

The strong effect in L5 IT neurons motivated a direct layer comparison. We imaged four Cux2-CreERT2 $\times$ Ai148 mice expressing GCaMP6 in layer 2/3 before and after clozapine, following the same protocol and analysis. These recordings yielded 137 short-range and 127 long-range pairs (Figure~\ref{fig:pair_counts}; Table~\ref{tab:stats}). Clozapine still produced a reduction in correlations in the L2/3 cohort ($^{*}$: $p<0.05$; $^{**}$: $p<0.01$; n.s.\ where indicated), but the reduction was visibly shallower than the L5 IT reduction when the 50\% contour lines from the pre-drug condition were overlaid on the post-drug heatmaps.

Rank-sum comparison of the drug-induced changes between the two cell-type cohorts confirmed this impression quantitatively. Cux2-CreERT2 $\times$ Ai148 versus Tlx3-Cre $\times$ Ai148 yielded short-range $p<0.005$ and long-range $p<10^{-8}$ (Figure~\ref{fig:pvalues}; Table~\ref{tab:stats}). The long-range bound was again several orders of magnitude more stringent than the short-range bound, reinforcing the interpretation that the antipsychotic effect targets inter-areal coupling preferentially, and does so in a manner that distinguishes superficial from deep cortical neurons. Together with the pan-neuronal baseline, these results place L5 IT neurons as the most strongly affected cortical population, followed by layer 2/3 neurons, with the pan-neuronal signal occupying an intermediate position that averages cell-type-specific effects.

\subsection{The psychostimulant amphetamine does not recapitulate antipsychotic decorrelation}

Antipsychotic drugs and the psychostimulant amphetamine both engage dopaminergic systems, but they produce opposite behavioural effects in humans and animal models~\citep{howes2009dopamine}. To test whether dopaminergic engagement alone is sufficient to produce L5 IT decorrelation, we imaged three Tlx3-Cre $\times$ Ai148 mice before and after a single intraperitoneal injection of amphetamine (4~$\mu$g/g body weight). Heatmaps and contour lines of correlation versus distance after amphetamine did not display the systematic reduction seen with the antipsychotics; the 50\% contour was not contracted in the way it was contracted after clozapine, aripiprazole, or haloperidol. Saline control injections (0.9\% NaCl, undiluted) likewise did not produce systematic changes in correlation structure.

This dissociation is informative. It argues against the simplest account in which any perturbation of the dopaminergic system uniformly alters L5 IT coupling, and instead supports a more specific account in which the circuit-level consequences of antipsychotic action differ qualitatively from those of a psychostimulant with overlapping dopaminergic pharmacology. Consistent with this interpretation, amphetamine exacerbates psychotic-like phenotypes in animal and human studies~\citep{howes2009dopamine, kambeitz2014alterations}, whereas the antipsychotics tested here reduce inter-areal coupling in exactly the cell class that is positioned to broadcast such coupling. The amphetamine control therefore localises the antipsychotic effect not at the level of receptor engagement per se, but at the level of the downstream circuit response.

\subsection{eGFP controls confirm that decorrelation is not a hemodynamic artifact}

A central concern for any widefield fluorescence measurement is that apparent neural-activity changes can be produced by drug-induced changes in cerebral blood flow and volume, which modulate background absorption and scattering. To control for this, we imaged 6 Tlx3-Cre mice that had been retro-orbitally injected with AAV-PHP.eB-Ef1$\alpha$-DIO-eGFP, yielding eGFP expression in L5 IT neurons, before and after a single injection of clozapine. We attempted an on-line hemodynamic correction using 405~nm isosbestic illumination of GCaMP6, but this correction could not fully recover the hemodynamic response measured in the eGFP recordings, indicating that accurate correction for hemodynamic response would require calibration on eGFP measurements rather than assumptions about the isosbestic wavelength alone.

Crucially, the drug-induced change in correlation-distance structure in the eGFP signal did not mirror the GCaMP6 decorrelation observed in Tlx3-Cre $\times$ Ai148 mice. This dissociation argues that the L5 IT decorrelation we observe under clozapine reflects neural activity rather than vascular contamination. In combination with the dose range explored for clozapine (0.2--10~$\mu$g/g body weight) and the convergent results with aripiprazole and haloperidol, the eGFP control makes a hemodynamic origin for the reported effect unlikely.

\subsection{Long-range interactions are preferentially decorrelated}

Collating the results across all antipsychotic comparisons, a consistent pattern emerges: long-range correlation reductions are more pronounced, and more significantly so, than short-range reductions. The pan-neuronal C57BL/6 comparison drew on 204 short-range and 192 long-range pairs; the Tlx3-Cre $\times$ Ai148 clozapine comparison on 182 and 148; the Tlx3 haloperidol/aripiprazole comparisons each on 114 short-range and 84 long-range pairs; and the Cux2 clozapine comparison on 137 and 127 pairs (Figure~\ref{fig:pair_counts}; Table~\ref{tab:stats}). Across these comparisons, the most stringent statistical bounds are consistently attached to long-range pairs: $p<10^{-5}$ for C57BL/6 versus Tlx3 and $p<10^{-8}$ for Cux2 versus Tlx3 (long-range), compared with $p<0.05$ and $p<0.005$ respectively for short-range (Figure~\ref{fig:pvalues}).

Taken together, the results support a unified picture. Antipsychotic drugs act most strongly on the correlated activity of L5 IT pyramidal neurons, in a manner that preferentially affects long-range inter-areal pairs, that is not reproduced by the psychostimulant amphetamine, that is not explained by hemodynamic contamination, and that is substantially weaker in layer 2/3 neurons. These observations are consistent with L5 IT neurons serving as a cell-type-specific cortical substrate through which multiple antipsychotic drugs, despite diverging in their receptor profiles, converge on a common circuit-level effect.

\section{Tables}

\begin{table}[htbp]
\centering
\caption{Mouse cohorts and cell-type labels used in the study. Counts correspond to the onset-response analyses of locomotion and visual flow. Clozapine, aripiprazole, haloperidol, amphetamine, and saline experiments used subsets of these cohorts as detailed in the Methods.}
\label{tab:cohorts}
\begin{tabular}{llllc}
\toprule
Mouse line & Indicator & Delivery / driver & Cortical population & N mice \\
\midrule
C57BL/6 & GCaMP6 & AAV-PHP.eB retro-orbital & Pan-neuronal & 6 \\
Emx1-Cre & GCaMP6 & AAV-PHP.eB retro-orbital (DIO) & Cortical excitatory & 4 \\
Scnn1a-Cre & GCaMP6 & Transgenic / AAV & Layer 4 (granular) & 7 \\
Tlx3-Cre $\times$ Ai148 & GCaMP6 & Ai148 reporter & L5 IT pyramidal & 14 \\
Ntsr1-Cre & GCaMP6 & Transgenic / AAV & L6 corticothalamic & 3 \\
PV-Cre & GCaMP6 & AAV-PHP.eB (DIO) & Parvalbumin interneurons & 2 \\
VIP-Cre & GCaMP6 & AAV-PHP.eB (DIO) & VIP interneurons & 6 \\
Sst-Cre & GCaMP6 & AAV-PHP.eB (DIO) & Somatostatin interneurons & 6 \\
Cux2-CreERT2 $\times$ Ai148 & GCaMP6 & Ai148 reporter (tamoxifen) & Layer 2/3 pyramidal & 3 \\
\bottomrule
\end{tabular}
\end{table}

\begin{table}[htbp]
\centering
\caption{Pharmacological agents, vendors, doses, and administration route. All agents were administered by a single intraperitoneal (i.p.) injection. Dose is expressed per gram body weight, except saline which was administered undiluted.}
\label{tab:drugs}
\begin{tabular}{llll}
\toprule
Drug & Vendor & Dose (per g body weight) & Vehicle \\
\midrule
Clozapine & Novartis & 0.2--10~$\mu$g/g & --- \\
Aripiprazole & Otsuka Pharmaceutical & 0.2~$\mu$g/g & --- \\
Haloperidol & Janssen & 0.1~$\mu$g/g & --- \\
Amphetamine & H\"anseler & 4~$\mu$g/g & --- \\
Saline (control) & --- & undiluted & 0.9\% NaCl in H$_2$O \\
\bottomrule
\end{tabular}
\end{table}

\begin{table}[htbp]
\centering
\caption{Pair counts and statistical bounds for the main correlation comparisons. All $p$-values are reported as upper bounds; rank-sum indicates a rank-sum test with family-wise error correction across $m=9$ genotype groups. C57BL/6 vs Tlx3-Cre $\times$ Ai148 and Cux2-CreERT2 $\times$ Ai148 vs Tlx3-Cre $\times$ Ai148 were the two principal cross-genotype contrasts.}
\label{tab:stats}
\begin{tabular}{lcccc}
\toprule
Comparison & Short-range pairs & Long-range pairs & Short-range $p$ & Long-range $p$ \\
\midrule
Pan-neuronal (C57BL/6) + clozapine & 204 & 192 & $<$0.001 / n.s. & $<$0.001 / n.s. \\
Tlx3-Cre $\times$ Ai148 + clozapine & 182 & 148 & $<$0.001 & $<$0.001 \\
Tlx3-Cre $\times$ Ai148 + haloperidol/aripiprazole & 114 & 84 & $<$0.05 & $<$0.001 \\
Cux2-CreERT2 $\times$ Ai148 + clozapine & 137 & 127 & $<$0.05 & $<$0.01 \\
\midrule
C57BL/6 vs Tlx3-Cre $\times$ Ai148 (rank-sum) & --- & --- & $<$0.05 & $<$10$^{-5}$ \\
Cux2-CreERT2 $\times$ Ai148 vs Tlx3-Cre $\times$ Ai148 (rank-sum) & --- & --- & $<$0.005 & $<$10$^{-8}$ \\
\bottomrule
\end{tabular}
\end{table}

\section{Discussion}

Our results identify layer 5 intratelencephalic pyramidal neurons as a cell-type-specific cortical substrate of antipsychotic drug action. Across three antipsychotics with distinct receptor profiles---clozapine, aripiprazole, and haloperidol---we observed a consistent reduction of inter-areal correlations that was most pronounced in L5 IT neurons, substantially weaker in layer 2/3 Cux2 neurons, and intermediate in pan-neuronal recordings. The reduction preferentially affected long-range pairs of cortical regions, with rank-sum $p$-value bounds reaching $10^{-5}$ for the C57BL/6 versus Tlx3 contrast and $10^{-8}$ for the Cux2 versus Tlx3 contrast on long-range pairs. Amphetamine, a psychostimulant that shares dopaminergic engagement with the antipsychotics but that has opposing clinical effects, failed to reproduce this decorrelation, and eGFP controls argued against a hemodynamic origin. Together, these observations connect a specific genetically defined cortical cell class to the acute circuit signature of antipsychotic action.

These findings are naturally framed within predictive-processing accounts of cortical function. L5 IT neurons are densely cortico-cortical and thus well-positioned to carry or gate top-down signals---such as the motor-based prediction of visual flow described by~\citet{leinweber2017sensorimotor} or the mismatch signals characterised by~\citet{keller2012sensorimotor} and~\citet{zmarz2016mismatch}---across dorsal cortex. In this view, antipsychotic drugs may dampen the broadcasting of internally generated predictions or correlated states, which could plausibly correspond to the ``calming'' of aberrant inter-areal coupling that has been proposed as a substrate of positive-symptom relief~\citep{uhlhaas2015oscillations, uhlhaas2013highfrequency, bahner2017hippocampal}. The observation that L2/3 Cux2 neurons are relatively spared is consistent with the idea that antipsychotic action spares local sensory computations and preferentially rebalances distributed signalling~\citep{harris2019hierarchical, siegle2021survey, keller2018predictive}.

The convergence of clozapine, aripiprazole, and haloperidol---despite their well-known differences in D2 and 5-HT2A affinities---on a common L5 IT decorrelation pattern speaks to the gap between receptor-level and circuit-level accounts of therapeutic action~\citep{howes2009dopamine, kambeitz2014alterations}. Receptor-centric models have been powerful for drug discovery but offer limited guidance for predicting inter-areal consequences. Our data suggest that the common factor across these antipsychotics, at the level of cortical dynamics, is a selective reduction of long-range coupling in cortico-cortically projecting deep-layer neurons. The dissociation from amphetamine further argues that not every pharmacological perturbation that engages dopamine produces this circuit-level signature; rather, the signature appears specifically tied to antipsychotic action.

Several features of our genetic and imaging strategy are worth emphasising. The use of Tlx3-Cre $\times$ Ai148 and Cux2-CreERT2 $\times$ Ai148 crosses, together with retro-orbital AAV-PHP.eB delivery in C57BL/6 mice, allowed the same imaging and analysis pipeline to be applied across pan-neuronal, L5 IT, and L2/3 populations, with only genetic driver identity varying across cohorts. Cre-dependent reporter lines such as Ai148 provide consistent, copy-controlled GCaMP6 expression across animals~\citep{daigle2018suite, chen2013ultrasensitive, bounds2023alloptical}, and AAV-PHP.eB enables high-efficiency noninvasive transduction for pan-neuronal comparisons~\citep{chan2017engineered}. Statistical analysis relied on hierarchical bootstrapping of time-course data~\citep{saravanan2020hierarchical} and family-wise error-corrected rank-sum tests for pairwise correlations, providing a conservative framework for cross-cohort comparisons. Anatomical context was provided by prior tracing and connectivity atlases that place Tlx3-positive neurons within the dense cortico-cortical projection network of dorsal cortex~\citep{liu2024celltype, harris2019hierarchical, matho2021genetic, biccn2021multimodal}.

Several limitations deserve explicit mention. First, widefield imaging samples predominantly superficial cortex: given a signal attenuation half-depth of about 100~$\mu$m, most fluorescence signal is expected to arise from the superficial ${\sim}100~\mu$m of tissue, even when the genetic driver labels deeper-layer neurons whose apical dendrites extend into these upper layers. Second, absolute cell-count comparisons across genotypes are not possible with widefield imaging; our conclusions rest on within-cohort drug effects and on their cross-cohort comparison, not on matching baseline activity levels. Third, the dose range explored for clozapine spans two orders of magnitude (0.2--10~$\mu$g/g), and we have not performed a formal dose-response analysis per genotype; the converging results nonetheless suggest that the effect is robust across this range. Fourth, the Cux2-CreERT2 cohort is smaller (3 mice) than the Tlx3-Cre cohort (14 mice), and the cross-genotype statistical test, while based on rank-sum comparison of pair-level correlation changes, is subject to the usual caveats of across-cohort comparisons. Finally, hemodynamic correction with 405~nm isosbestic illumination did not fully recover the response measured in eGFP controls, and we rely on the eGFP experiments themselves as the definitive control.

A further constraint is that our measurements are acute: each mouse received a single injection of one drug, and the analyses contrast pre- and post-injection data collected within the same session. Our design therefore does not address chronic effects, which are the clinically relevant regime for most antipsychotic prescriptions, nor does it address the relationship between acute cortical decorrelation and behavioural outcomes such as symptom relief. Nevertheless, the reproducibility of the L5 IT decorrelation across three antipsychotics, its absence under amphetamine, its weakness in L2/3 Cux2 neurons, and its confirmation against eGFP controls provide a robust acute signature that can anchor future longitudinal studies.

Looking forward, the results motivate several directions. First, cell-type- and layer-targeted manipulations of L5 IT neurons---chemogenetic silencing, optogenetic perturbation, or projection-specific interventions informed by anatomical atlases~\citep{liu2024celltype, harris2019hierarchical, matho2021genetic}---can test whether modulating L5 IT coupling causally mediates behavioural phenotypes relevant to psychosis-like states in mice. Second, pairing widefield imaging with electrophysiological recordings will link the correlation-level decorrelation reported here to spike-level changes in inter-areal synchrony, which is the level at which oscillatory abnormalities have been reported in human psychotic disorders~\citep{uhlhaas2015oscillations, uhlhaas2013highfrequency}. Third, from a translational perspective, the convergence of structurally diverse antipsychotics on a single cell-type-specific cortical effect suggests that next-generation compounds could be screened, at least in part, by their ability to reproduce L5 IT decorrelation in similar preparations, providing a circuit-level complement to receptor-based screening pipelines. More broadly, by tying a classical clinical pharmacology to a molecularly defined cortical cell class, the present study contributes a bridge between the receptor-centred and circuit-centred traditions in the neuroscience of antipsychotic action.

\bibliographystyle{unsrtnat}
\bibliography{references}

\end{document}
